%----------------------------------------------------------------------------------------
%	DESCRIPTION OF THE INTERACTION BETWEEN THE STUDENT AND THE INTERNAL TUTOR
%----------------------------------------------------------------------------------------
\section*{Interazione tra studente e tutor aziendale}
% Personalizzare definendo le modalità di interazione col tutor aziendale

Lo studente avrà un'interazione diretta e quotidiana con il tutor aziendale \nomeTutorAziendale\ \cognomeTutorAziendale\ (\ruoloTutorAziendale\ di \ragioneSocAzienda), che lo seguirà personalmente per l'intera durata del tirocinio.

I canali di comunicazione utilizzati sono i seguenti:
\begin{itemize}
    \item \textbf{Posta elettronica}: tramite l'indirizzo aziendale del tutor (\emailTutorAziendale) e l'indirizzo dello studente (\emailStudente), per comunicazioni formali, condivisione di documentazione e materiale di studio;
    \item \textbf{WhatsApp}: tramite il numero del tutor (\telTutorAziendale), per comunicazioni rapide, chiarimenti operativi e coordinamento giornaliero;
\end{itemize}

A complemento dell'interazione diretta, le attività di code review verranno svolte sulla piattaforma GitLab, dove lo studente proporrà le proprie modifiche tramite merge request che il tutor analizzerà e commenterà, fornendo feedback puntuali sul codice prodotto.

Periodicamente, inoltre, sono previsti momenti di confronto strutturato con il tutor e con eventuali ulteriori stakeholders aziendali, finalizzati a verificare lo stato di avanzamento delle attività, chiarire eventuali dubbi sugli obiettivi e, se necessario, aggiornare il piano di lavoro.